\section{Related Work}

\paragraph{Interpretation of Code-Generation Models.} Post-hoc methods associate input features, concepts, or rationales with model behavior~\cite{sundararajan2017integrated,kim2018tcav,lundberg2017shap,deyoung2020eraser}. For code generation, CodeQ~\cite{khati2026codeq} maps token-level rationales to programming concepts, while NLPerturbator~\cite{chen2026nlperturbator} probes robustness under natural-language variations. These methods improve inspectability but do not construct task-scoped security requirement candidates or evaluate their effects on independently measured security outcomes.

\paragraph{Causal Analysis of Code Models.} Causal studies examine how specified factors affect code-model behavior. Palacio et al.~\cite{palacio2024causation} study software and model interventions on neural code-model predictions, while Galeras~\cite{rodriguezcardenas2023benchmarking} provides causal-evaluation testbeds for generative software tasks and predefined prompt-engineering treatments. Ji et al.~\cite{ji2025causality} model prompt--code concept relations, and CAM~\cite{lyu2026cam} attributes multi-agent outcomes to intermediate traces. More directly related to prompt intervention, Velasco et al.~\cite{velasco2026smells} analyze how generation strategy, model characteristics, and prompt formulation affect code-smell propensity and evaluate prompt-based mitigation. Causal analyses have also covered privacy leakage through training dynamics and membership inference under semantics-preserving code transformations~\cite{yang2025privacy,yang2025membership}. These works analyze specified code, model, training, trace, or prompt factors. \model also uses a fixed requirement catalog, but instantiates context-scoped candidates from individual tasks, prioritizes them under a shared budget, and freezes matched prompt policies before randomized held-out evaluation. The Prompt TSG constrains candidate semantics rather than encoding causal relations; observational discovery selects what to test, while randomization estimates policy effects on separate security and functionality outcomes.

\paragraph{Hardening LLMs for Secure Code.} Prompt-level studies have examined security reminders, CWE-specific instructions, and iterative repair~\cite{tony2025prompting,bruni2025promptbenchmark}. Other approaches use instruction tuning (SafeCoder), retrieved demonstrations (SecCoder), continuous steering (SVEN), prompt optimization (PromSec), constrained decoding (CodeGuard+), or representation intervention (Thea)~\cite{he2024safecoder,zhang2024seccoder,he2023sven,nazzal2024promsec,fu2024codeguard,elhusseini2026thea}. \model instead determines which catalog-bounded requirements apply to a task, which candidates warrant evaluation, and whether their ADD or REMOVE policies change security, with functionality and requirement specificity measured separately.

\paragraph{Evaluation Benchmarks for Secure Code Generation.} SecurityEval~\cite{siddiq2022securityeval} and LLMSecEval~\cite{tony2023llmseceval} established CWE-indexed prompt suites primarily evaluated through static analysis, while SALLM~\cite{siddiq2024sallm}, CWEval~\cite{peng2025cweval}, BaxBench~\cite{vero2025baxbench}, and SeCodePLT~\cite{nie2025secodeplt} incorporate executable checks or dynamic exploits. Because apparent security improvements can reflect functional degradation or measurement failures~\cite{dai2026rethinking}, \model keeps code validity, Oracle support, oracle-evaluable security, functionality, and secure-and-functional joint success as separate evaluation outcomes.
